\section{Evaluation Protocol}\label{sec:experiments}

\subsection{Development subset and full-scale protocol}

The development corpus contains 150 clips and all 128 body configurations,
for 19,200 variants. These clips are reserved within the full training split
and are not a held-out generalization benchmark. Small-scale experiments
investigate architecture and reference-processing choices with a fixed
simulation budget. The full-scale run uses the 18,855 training clips in
Table~\ref{tab:dataset_counts}, with validation membership fixed throughout.

The selected full-scale experiment is slot/type attention trained directly on
refined references. Its code retains a historical \texttt{stage2} name, but
this run does not use neutral pretraining followed by shape transfer.
Table~\ref{tab:training_parameters} records the shared optimization settings
and the actual environment/batch sizes used after the initial memory-related
launch adjustments. Run-specific resolved configurations take precedence
over older launch examples.

\begin{table}[t]
\centering
\small
\caption{Training configuration and budgets. Environment and minibatch counts
are per GPU. The full-scale frame count is a recorded stopping snapshot,
not a budget matched to the development runs.}\label{tab:training_parameters}
\begin{tabularx}{\linewidth}{@{}lYY@{}}
\toprule
Setting & Corrected 150-clip runs & Full-scale selected run \\
\midrule
Parallel GPUs per policy & 1 & 6 \\
Environments per GPU & 4,096 & 2,048 \\
PPO minibatch size per GPU & 16,384 & 8,192 \\
Rollout length & 32 control steps & 32 control steps \\
Mini-epochs per update & 1 & 1 \\
Actor / critic learning rates &
$2\times10^{-5}$ / $10^{-4}$ & $2\times10^{-5}$ / $10^{-4}$ \\
Discount / GAE / ratio clip & 0.99 / 0.95 / 0.2 & 0.99 / 0.95 / 0.2 \\
Gradient norm clip & 50 & 50 \\
Evaluation cadence & 200 epochs & 256 epochs \\
Training seed & 0--2 (refined A--D); 0 (unrefined E--H) & Recorded in saved run configuration \\
Reference residency & All 150 clips & 256, later 128 clips per GPU \\
Training frames & 786,432,000 per run (6,000 epochs) & Approximately $11.07\times10^9$ recorded \\
Status of reported evidence & Completed, Table~\ref{tab:architecture_seeds} & Epoch-28,170 checkpoint, Table~\ref{tab:test_results} \\
\bottomrule
\end{tabularx}
\end{table}

\subsection{Design comparisons and aggregation}\label{sec:ablation_design}
We compare four architectures: A, temporal MLP; B, temporal attention; C,
attention with slot/type embeddings; and D, C with actor-only adaLN-Zero.
All retain morphology inputs in actor and critic. Refined-reference runs
use three seeds (0, 1, 2), a matched 6,000-epoch budget, and the same cadence.
Each run is summarized by its last eight evaluations (epochs 4,600--6,000);
architecture means and sample standard deviations are then computed across
three run summaries. Within-run evaluations are not training replicates.

The full-scale model predates these supporting reruns. They reassess its
design rather than establish prospective selection or statistical superiority.
Models differ in parameter count; compute details are in Appendix~\ref{app:history}.
The original seed-0 grid is retained in Table~\ref{tab:ablation_results}.
Unrefined cases E--H have one seed each. Refined and unrefined policies track
their own targets, preventing a causal policy-level refinement conclusion.
The paired offline reference audit provides the evidence for artifact reduction.
\begin{table}[t]
\centering
\small
\caption{Corrected 150-clip grid: four architectures $\times$ two reference
versions, seed 0. Refined A--D also have seed-1/2 repeats in Table~\ref{tab:architecture_seeds}.}\label{tab:ablations}
\begin{tabularx}{\linewidth}{@{}llYY@{}}
\toprule
Case & References & Architecture & Comparison purpose \\
\midrule
A & Refined & Wide temporal MLP; beta concatenation &
Temporal MLP baseline. \\
B & Refined & Basic temporal attention &
Attention versus flattened temporal input. \\
C & Refined & Attention with slot/type embeddings &
Explicit temporal/source identity. \\
D & Refined & C + actor-only adaLN-Zero &
Morphology modulation with critic unchanged. \\
E & Unrefined & As C &
Reference version at fixed architecture (slot/type). \\
F & Unrefined & As A &
Reference version at fixed architecture (temporal MLP). \\
G & Unrefined & As B &
Reference version at fixed architecture (attention). \\
H & Unrefined & As D &
Reference version at fixed architecture (slot/type + adaLN). \\
\bottomrule
\end{tabularx}
\end{table}
\subsection{Morphology-matched evaluation and checkpoint scope}\label{sec:evaluation_protocol}

Each reference is rolled out only in an environment using its exact asset
identity (Figure~\ref{fig:evaluation_protocol}). A selected reference
therefore cannot be assigned to a different body simply because an
environment slot is available. Multiple rollout batches may be needed when
several references require the same body.

Routine monitoring selects one variant per clip. With seed $s=42$, stable
clip identity $c$, panel index $p$, and $N_c$ available variants, the choice is
\begin{equation}
    k(c,p) =
    \left(
       \operatorname{uint64}_{\mathrm{big}}
       \bigl(\operatorname{SHA256}(s\!:\!c)_{0:8}\bigr)
       +p
    \right)\bmod N_c .
    \label{eq:shape_panel}
\end{equation}
Here $s\!:\!c$ is the UTF-8 string formed as \texttt{seed:clip\_id}.
For scheduled monitoring, $p$ is the integer quotient of the checkpointed
epoch and evaluation cadence. This produces deterministic staggered shape
choices that rotate across evaluations and remain reproducible after resume.
The validation \emph{clip list} is fixed; the selected \emph{shape panel}
changes. Thus successive scores are not repeated measurements on an
identical set of clip/shape pairs.

Evaluation uses the mean action, reference initialization, and no gradient
tracking. It disables training termination and records errors over the
valid reference horizon, capped at 600 control steps (20\,s at 30\,Hz).
Success therefore refers to the evaluated horizon, not necessarily the
entire duration of a longer source clip.

Final model comparisons should freeze the checkpoint, clip/shape panel,
reference version, simulator configuration, and seed. Validation is available
for model selection and has been repeatedly monitored; it is not an untouched
test set. The independent test manifest remains outside training and model
selection; its one-time evaluation is reported in
Section~\ref{sec:test_results}. The downloaded epoch-28,170
\texttt{last.ckpt} is distinct from the epoch-28,159 monitoring snapshot
in Table~\ref{tab:full_results}; the standalone validation and test results
therefore refer explicitly to the downloaded checkpoint.

\begin{figure}[t]
\centering
\begin{tikzpicture}[>=latex, node distance=0.35cm, every node/.style={font=\small}]
\node[draw, rounded corners, fill=blue!8, align=center, minimum width=2.05cm, minimum height=0.72cm] (clip) {clip identity\\+ manifest};
\node[draw, rounded corners, fill=green!10, right=of clip, align=center, minimum width=2.15cm, minimum height=0.72cm] (pick) {SHA-256\\shape pick};
\node[draw, rounded corners, fill=orange!12, right=of pick, align=center, minimum width=2.15cm, minimum height=0.72cm] (asset) {matching\\SMPLSim asset};
\node[draw, rounded corners, fill=purple!10, right=of asset, align=center, minimum width=1.9cm, minimum height=0.72cm] (roll) {mean-action\\rollout};
\node[draw, rounded corners, fill=gray!12, right=of roll, align=center, minimum width=1.75cm, minimum height=0.72cm] (metrics) {per-motion\\metrics};
\draw[->, thick] (clip) -- (pick);
\draw[->, thick] (pick) -- (asset);
\draw[->, thick] (asset) -- (roll);
\draw[->, thick] (roll) -- (metrics);
\node[below=0.32cm of pick, align=center] {one body per clip\\(rotating panel)};
\node[below=0.32cm of roll, align=center] {no gradients\\capped horizon};
\end{tikzpicture}
\caption{Standalone evaluation protocol. The clip list is fixed, while a
reproducible hash-and-panel rule selects one body variant per clip; the
selected reference and simulation asset therefore always match.}
\label{fig:evaluation_protocol}
\end{figure}

\subsection{Metrics and aggregation}\label{sec:metrics}

For each evaluated clip/shape pair $i$, define its per-frame mean body
position error
\begin{equation}
    d_{i,t} = \frac{1}{B}\sum_{j=1}^{B}
                 \|p_{i,t,j}-\hat{p}_{i,t,j}\|_2 .
\end{equation}
The configured success indicator is
\begin{equation}
    S_i = \mathbf{1}\!\left[
             \max_{t<T_i}d_{i,t}\leq 0.5\ \mathrm{m}
          \right],
    \label{eq:success}
\end{equation}
where $T_i$ is the evaluated horizon. This threshold is deliberately distinct
from the smoother training reward. Passing it does not imply fine-grained
pose accuracy or biomechanical realism.

We also report time-averaged mean body-position error, mean global
body-rotation geodesic error, and the time average of the largest body-position
error within each frame. The last quantity is the logged
\emph{mean maximum-joint error}; it is not the maximum over the entire dataset.
Global body-rotation error is not a local joint-angle/DOF error.

Within an evaluator rank, errors are averaged over valid frames per motion
and then over motions, including failed rollouts. The distributed training
logger aggregates scalar summaries across ranks; the monitoring table
retains those logged values rather than reconstructing exact pooled counts.
Final offline comparisons should aggregate saved per-motion records over
unique clip/shape pairs, avoiding unequal weighting of partially filled
batches or duplicate padded rollouts.

Trajectory smoothness uses the implementation's windowed normalized-jerk
score, with the definition and numerical regularizers in
Appendix~\ref{app:smoothness}. Lower scores indicate lower jerk under that
fixed convention. The accompanying high-jerk statistic is the percentage
of windows in which any body exceeds threshold 6,500; the logger calls it
a frame percentage, but it is window-based.


\subsection{Research questions, populations, and checkpoint scope}
Table~\ref{tab:evidence_populations} maps the evidence to the three questions.
Held-out identity is clip-grouped, not a demonstrated subject/recording split.
\begin{table}[t]
\centering\small
\caption{Evaluation populations and their inference scope. Each unseen panel
has 128 bodies; the 233-clip and 23-clip subsets partition its 256 clips.}
\label{tab:evidence_populations}
\begin{tabularx}{\linewidth}{@{}lYYr@{}}
\toprule
Question & Motion exposure & Body exposure & Pairs \\
\midrule
RQ1: validation / test & Held-out clips & Trained; one per clip & 1,048 each \\
RQ1: development replay & 150 training clips & All 128 trained & 19,200 \\
RQ2: inside / outside & 233 training clips & All 128 unseen per panel & 29,824 each \\
RQ2: diagnostic subset & 23 val/test clips; seen by full-data & Same unseen panels & 2,944 each \\
RQ2: trained control & 150 training clips & All 128 trained & 19,200 \\
RQ3: design comparison & Development clips & Trained bodies & Three refined seeds \\
\bottomrule
\end{tabularx}
\end{table}

\paragraph{Unseen-body construction and replay.}
Each panel contains 64 new coefficient vectors instantiated on male and female
SMPL templates. Inside vectors lie in $[-3,3]^{10}$; outside vectors have
$\|\beta\|_\infty\in(3,5]$, with 16 vectors (32 bodies) in each of four
half-unit tiers. Inside denotes membership of the training coefficient box,
not demonstrated membership of the convex hull of trained bodies. Both
panels use matching HUMOS references and assets, frame-zero offsets and
full reference refinement. The build report records finite complete clip/body
panels and zero penetration under its refinement diagnostic; it does not
certify dynamic feasibility.

All nine runs completed on an A40 on 9 October 2026. Replay uses mean actions,
reference initialization, 32 environments per asset (4,096 environments),
and at most 210 control steps (7\,s). Success follows Eq.~\eqref{eq:success}.
The 233 training-split identities isolate the body axis; the other 23
validation/test identities remain separate. The trained-body control uses a
different 150-clip panel and refinement run, so only within-dataset checkpoint
deltas are paired comparisons.

\paragraph{Checkpoint and uncertainty controls.}
The original \texttt{refined} checkpoint is epoch 28,170; the prespecified
\texttt{fulldata} candidate is epoch 34,000 and trained on all original motion
splits. Both are frozen for replay. Duration and motion coverage change
together, preventing causal attribution. Each dataset has two original-checkpoint
runs and one candidate run. Paired 95\% bootstrap intervals resample clips,
retaining the full fixed body panel; they do not resample training seeds or
unseen bodies. The report's ``beyond noise'' rule requires an interval excluding
zero and a candidate delta larger in magnitude than the original-checkpoint
rerun delta. With no checkpoint ladder, that comparison measures simulator
variability only. Per-body regression intervals are exploratory and unadjusted
for multiple testing.

On all 256 clips, identical-checkpoint success labels agree on 99.33\% of
inside and 98.68\% of outside pairs, versus 99.96\% in the control. Pooled
success differs by at most 0.05 pp, but individual pairs can change by metres
in peak error. Small pooled changes therefore coexist with local instability.
Checkpoint hashes and the retained report are documented in
Appendix~\ref{app:provenance}; complete per-body tables are in
Appendix~\ref{app:unseen_tables}.
